\subsection{Correlated queries across folded oracles}\label{sec:kbfoldfold:joint}

The polynomials $Q_0,\ldots,Q_m$ share the \emph{same} random mask $R$. Consequently their oracle responses are correlated. We now give the exact joint distribution for any fixed set of observations. This extends the single-oracle masking criterion of Chapter~4, but it is not yet an online simulation theorem for a malicious adaptive verifier.

\begin{definition}[Fixed fold-query schedule]\label{def:kbff:schedule}
Retain the fields and domains of Section~\ref{sec:kbfoldfold:codewords}.
Fix a complete challenge vector $r\in\F^m$ and an integer $k\geq0$.
A \emph{fixed query schedule} is a sequence
\[
 \Sigma=((j_1,\xi_1),\ldots,(j_k,\xi_k)),\qquad
 0\leq j_h\leq m,\quad \xi_h\in\Omega_{j_h}.
\]
The schedule is fixed independently of the mask $R$; it may depend on
the already fixed challenge vector $r$.
\end{definition}

\begin{definition}[Witness and mask observation maps]\label{def:kbff:observation-maps}
For such a fixed $r$ and $\Sigma$, define the maps
\[
 \mathcal S_{r,\Sigma}:\F^{I_m}\to\F^{k+1},
 \qquad
 \mathcal T_{r,\Sigma}:\F[X]_{<B}\to\F^{k+1}
\]
by
\[
 \mathcal S_{r,\Sigma}(f)
 :=\left(U_{f_{j_h}}^{(m-j_h)}(\xi_h)\right)_{h=1}^k
       \mathbin{\|}(0),
\]
\[
 \mathcal T_{r,\Sigma}(R)
 :=\left(\xi_h^{\,2^{m-j_h}}R_{j_h}(\xi_h)\right)_{h=1}^k
       \mathbin{\|}(\rho_r(R)),
\]
where $\mathbin{\|}$ denotes concatenation with one final coordinate,
$f_j$ and $R_j$ are as in Definition~\ref{def:kbff:masked-chain},
and $\rho_r(R)$ is the scalar in Theorem~\ref{thm:kbff:linear-closure}.
\end{definition}

\begin{lemma}[Both observation maps are linear]\label{lem:kbff:observation-linear}
The maps $\mathcal S_{r,\Sigma}$ and $\mathcal T_{r,\Sigma}$
are $\F$-linear.
\end{lemma}
\begin{proof}
The map $f\mapsto f_j$ is linear by Lemma~\ref{lem:kbff:prefix-linear};
$f_j\mapsto U_{f_j}^{(m-j)}$ is linear by the finite-sum definition;
evaluation at a fixed $\xi_h$ is linear. Thus each of the first
$k$ coordinates of $\mathcal S_{r,\Sigma}$ is linear, and the final
zero coordinate is linear.
For $\mathcal T_{r,\Sigma}$, each fold map $\pfold_{r_i}$ is linear,
so $R\mapsto R_{j_h}$ is linear. Evaluation at $\xi_h$ and
multiplication by the fixed scalar
$\xi_h^{2^{m-j_h}}$ are linear. Finally
$R\mapsto\rho_r(R)=\mathsf{PFold}^{[m]}_r(R)$ is linear by the same
fold-linearity argument. Hence all coordinates are linear.
\end{proof}

\begin{lemma}[Exact observation decomposition]\label{lem:kbff:observation-decompose}
Let $Q_j$ be the masked fold polynomials. Define the observed vector
\[
 \mathcal O_{r,\Sigma}(f,R)
 :=\left(Q_{j_h}(\xi_h)\right)_{h=1}^k
     \mathbin{\|}\bigl(\rho_r(R)\bigr).
\]
Then
\[
 \boxed{
 \mathcal O_{r,\Sigma}(f,R)
  =\mathcal S_{r,\Sigma}(f)+\mathcal T_{r,\Sigma}(R).
 }
\]
\end{lemma}
\begin{proof}
For each query $(j_h,\xi_h)$, Theorem~\ref{thm:kbff:masked-invariant}
gives
\[
 Q_{j_h}(\xi_h)
 =U_{f_{j_h}}^{(m-j_h)}(\xi_h)
   +\xi_h^{2^{m-j_h}}R_{j_h}(\xi_h).
\]
These are exactly the $h$-th coordinates of
$\mathcal S_{r,\Sigma}(f)$ and $\mathcal T_{r,\Sigma}(R)$.
The last coordinate of the first vector is $0$, and the last
coordinate of the second is $\rho_r(R)$, which also equals the
last coordinate of $\mathcal O_{r,\Sigma}(f,R)$.
Thus all $k+1$ coordinates agree.
\end{proof}

\begin{theorem}[Exact joint fixed-schedule masking criterion]\label{thm:kbff:joint-mask}
Assume $\F$ is finite and $R$ is uniform on the vector space
$\F[X]_{<B}$. For fixed $r$ and $\Sigma$, the vector
$\mathcal O_{r,\Sigma}(f,R)$ is uniform on the affine coset
\[
 \mathcal S_{r,\Sigma}(f)
   +\operatorname{im}\mathcal T_{r,\Sigma}
 \subseteq\F^{k+1}.
\]
For any two tables $f,g:I_m\to\F$,
\[
 \boxed{
 \Delta\left(
 \mathcal O_{r,\Sigma}(f,R),\mathcal O_{r,\Sigma}(g,R)
 \right)
 =
 \begin{cases}
 0,&\mathcal S_{r,\Sigma}(f-g)
        \in\operatorname{im}\mathcal T_{r,\Sigma},\\
 1,&\text{otherwise}.
 \end{cases}}
\]
\end{theorem}
\begin{proof}
By Lemma~\ref{lem:kbff:observation-linear},
$\mathcal T_{r,\Sigma}$ is linear. Lemma~\ref{lem:kbff:observation-decompose}
expresses the random vector as the sum of a fixed vector
$\mathcal S_{r,\Sigma}(f)$ and the image of a uniform random
element of the mask space. Lemma~\ref{lem:security:affine-uniform}
therefore gives the uniform coset distribution.
Apply Lemma~\ref{lem:security:coset-criterion} with
$V=\F^{k+1}$, $H=\operatorname{im}\mathcal T_{r,\Sigma}$,
and secret vectors $\mathcal S_{r,\Sigma}(f)$ and
$\mathcal S_{r,\Sigma}(g)$. Their difference equals
$\mathcal S_{r,\Sigma}(f-g)$ by linearity, proving both cases.
\end{proof}

\begin{corollary}[Sufficient condition for witness-independent fixed-schedule answers]\label{cor:kbff:joint-surjective}
If $\mathcal T_{r,\Sigma}$ is surjective onto $\F^{k+1}$,
the joint observed vector is uniform on $\F^{k+1}$ for every
witness $f$, hence independent of $f$.
\end{corollary}
\begin{proof}
Surjectivity gives $\operatorname{im}\mathcal T_{r,\Sigma}
=\F^{k+1}$. Its unique affine coset is the whole target space.
Apply the uniform-coset conclusion of Theorem~\ref{thm:kbff:joint-mask}.
\end{proof}

\begin{lemma}[Dimension obstruction to surjectivity]\label{lem:kbff:rank-obstruction}
If $\mathcal T_{r,\Sigma}$ is surjective, then
\[
 k+1\leq B.
\]
\end{lemma}
\begin{proof}
The source $\F[X]_{<B}$ has dimension $B$ over $\F$.
The target $\F^{k+1}$ has dimension $k+1$.
The image of a linear map has dimension at most that of its source.
Surjectivity would give
$k+1=\dim\operatorname{im}\mathcal T_{r,\Sigma}\leq B$.
\end{proof}

\begin{remark}[Final line coefficients are correlated with earlier openings]
The last observation $\rho_r(R)$ is a linear functional of the
\emph{same} mask that randomizes the earlier oracle values.
Even if each marginal query is uniform, the joint distribution
may distinguish two witnesses unless
$\mathcal S_{r,\Sigma}(f-g)$ lies in the mask image.
Theorem~\ref{thm:kbff:joint-mask} is therefore the exact
fixed-schedule proof obligation for a candidate simulator.
When a verifier chooses later queries after seeing earlier
answers, the schedule is no longer fixed; one must use the
conditional-coset description in
Lemma~\ref{lem:security:conditioned-mask} for every adaptive history.
\end{remark}

\begin{remark}[Not a complete zero-knowledge theorem]
The observations here do not include commitment roots, Merkle
authentication paths, the sumcheck round-polynomial messages, or
adaptively selected challenge/query schedules. No claim that
$\mathcal T_{r,\Sigma}$ is always surjective is made.
Consequently the chapter supplies a necessary-and-sufficient
linear-algebraic test for a \emph{specified} fixed schedule,
not zero knowledge for the complete KBFold protocol.
\end{remark}
